\chapter{Executive Summary}\label{ch:summary}

The Harp is a digital sympathetic-string pedal. Behind the guitar sit up to 24 digitally modelled strings that nobody plays. They ring when the guitar sounds a note they are tuned to, as the taraf strings of a sitar or the under-strings of a Hardanger fiddle do. The pedal listens to the chord being played and keeps retuning the strings to it, so every note leaves a halo in the player's own harmony. The intended sound is tuned ambience: it is not a reverb and not a shimmer.

The category neighbours are the Hologram Microcosm, the Chase Bliss Mood and the Strymon ambience pedals, at 300 to 450\,EUR. None of them does chord-following sympathetic strings.

\section{Hardware in One Paragraph}
The Harp is a cut-down Timekeeper (25-Z01-0001) on the hardware set shared with The Alchemist (26-A02-0001) and The Relic (26-A03-0003). It uses the Timekeeper's AK4621EF codec, W25Q128 QSPI flash, buck converter, input and output stages, and The Relic's relay bypass. The MCU is an STM32H750VBT6 in LQFP-100 rather than the Timekeeper's LQFP-144, and there is no SDRAM: the 24 strings need about 100\,KB of delay memory and the H750 has 1\,MB internal. There is no screen. A 0.91\,in OLED strip shows the tuning, and four pots and three ON-ON-ON toggles replace the Timekeeper's encoders. Everything sits on one four-layer 56 by 90\,mm board in a Hammond 125B, with the jacks, DC jack, knobs and footswitches at the same positions and heights as The Alchemist and The Relic.

\section{Status at Revision A}
\begin{itemize}
  \item \textbf{Schematic.} A first draft of seven sheets and 190 parts, generated by \texttt{scripts/harp\_schematic/build\_layout.py} (stereo input and OLED display added 2026-10-08) and net-traced against its specification. It has not yet been opened in KiCad, and ERC has not been run.
  \item \textbf{Board.} The Relic's outline, four-layer stack and design rules, with 28 face, jack and bypass footprints placed. The remaining parts are not placed and nothing is routed.
  \item \textbf{Firmware.} A requirements draft exists (\texttt{docs/firmware-requirements.md}); no code is written. The engine is to be developed first on the Timekeeper board.
  \item \textbf{Measurements.} None. The CPU load of the string bank, the chord tracker's accuracy and every audio figure in this report are estimates or targets.
\end{itemize}

\section{Decisions Taken}
Thirteen design decisions were recorded between 2026-10-04 and 2026-10-06; they are listed in Chapter~\ref{ch:architecture}. The ones that shape everything else are: one relay with true bypass on OUT~L only and a Trails mode in firmware for stereo players; one board unless placement fails; pots rather than encoders, with four knobs and three ON-ON-ON toggles laid out as on The Alchemist; an STM32H750VBT6 with no SDRAM; the dry signal through the codec; and an OPA2365 input buffer on the 5\,V rail with a jack-to-ADC gain of 0.59.

\section{What Comes Next}
Firmware comes first, on the Timekeeper board. It proves the sound and gives the CPU figure before any Harp PCB is committed. In parallel, the schematic is opened in KiCad for ERC, and placement decides whether the one-board plan holds. Chapter~\ref{ch:risks} lists the open items and the order of work.
